% Options for packages loaded elsewhere
\PassOptionsToPackage{unicode}{hyperref}
\PassOptionsToPackage{hyphens}{url}
\documentclass[
  11pt,
  letterpaper,
]{article}
\usepackage{xcolor}
\usepackage[margin=1in]{geometry}
\usepackage{amsmath,amssymb}
\setcounter{secnumdepth}{-\maxdimen} % remove section numbering
\usepackage{iftex}
\ifPDFTeX
  \usepackage[T1]{fontenc}
  \usepackage[utf8]{inputenc}
  \usepackage{textcomp} % provide euro and other symbols
\else % if luatex or xetex
  \usepackage{unicode-math} % this also loads fontspec
  \defaultfontfeatures{Scale=MatchLowercase}
  \defaultfontfeatures[\rmfamily]{Ligatures=TeX,Scale=1}
\fi
\usepackage{lmodern}
\ifPDFTeX\else
  % xetex/luatex font selection
\fi
% Use upquote if available, for straight quotes in verbatim environments
\IfFileExists{upquote.sty}{\usepackage{upquote}}{}
\IfFileExists{microtype.sty}{% use microtype if available
  \usepackage[]{microtype}
  \UseMicrotypeSet[protrusion]{basicmath} % disable protrusion for tt fonts
}{}
\makeatletter
\@ifundefined{KOMAClassName}{% if non-KOMA class
  \IfFileExists{parskip.sty}{%
    \usepackage{parskip}
  }{% else
    \setlength{\parindent}{0pt}
    \setlength{\parskip}{6pt plus 2pt minus 1pt}}
}{% if KOMA class
  \KOMAoptions{parskip=half}}
\makeatother
\usepackage{longtable,booktabs,array}
\usepackage{calc} % for calculating minipage widths
% Correct order of tables after \paragraph or \subparagraph
\usepackage{etoolbox}
\makeatletter
\patchcmd\longtable{\par}{\if@noskipsec\mbox{}\fi\par}{}{}
\makeatother
% Allow footnotes in longtable head/foot
\IfFileExists{footnotehyper.sty}{\usepackage{footnotehyper}}{\usepackage{footnote}}
\makesavenoteenv{longtable}
\setlength{\emergencystretch}{3em} % prevent overfull lines
\providecommand{\tightlist}{%
  \setlength{\itemsep}{0pt}\setlength{\parskip}{0pt}}
\usepackage{bookmark}
\IfFileExists{xurl.sty}{\usepackage{xurl}}{} % add URL line breaks if available
\urlstyle{same}
\hypersetup{
  hidelinks,
  pdfcreator={LaTeX via pandoc}}

\author{}
\date{}

\begin{document}

\section{The Qualitative Nature of One, Two, and Three: Structural Role
Assignment in Minimal Recursive
Systems}\label{the-qualitative-nature-of-one-two-and-three-structural-role-assignment-in-minimal-recursive-systems}

\textbf{Author:} J. W. Milton, Clarity Coalition

\textbf{Version:} 1.0

\textbf{Date:} 17 April 2026

\textbf{Keywords:} integers; triadic roles; N-D-C framework; structural
role assignment; philosophy of mathematics; recursion

\begin{center}\rule{0.5\linewidth}{0.5pt}\end{center}

\subsection{Abstract}\label{abstract}

The tholonic framework assigns three functionally distinct roles to the
variables of minimal recursive systems: a unity state \(N\)
(negotiation), a bounding parameter \(D\) (definition/limitation), and
an integrating parameter \(C\) (contribution). Prior work in this series
establishes that exactly three such roles are structurally necessary
(\href{https://github.com/baardev/truevalue/raw/main/docnav/Research/papers/3_minimal-recursive-triadic-framework/3_minimal-recursive-triadic-framework.pdf}{paper
3}) and that they recur across recurrences generating classical
constants
(\href{https://github.com/baardev/truevalue/raw/main/docnav/Research/papers/1_recursive-tholonic-five-constants/1_recursive-tholonic-five-constants.pdf}{paper
1}), supply-chain transparency scoring
(\href{https://github.com/baardev/truevalue/raw/main/docnav/Research/papers/2_supply-chain-transparency-tvpci/2_supply-chain-transparency-tvpci.pdf}{paper
2}), game-theoretic equilibria
(\href{https://github.com/baardev/truevalue/raw/main/docnav/Research/papers/4_game-theoretic-triadic-balance/4_game-theoretic-triadic-balance.pdf}{paper
4}), and twistor geometry
(\href{https://github.com/baardev/truevalue/raw/main/docnav/Research/papers/5_tholonic-twistor-connection/5_tholonic-twistor-connection.pdf}{paper
5}). What has not been argued formally is whether this role assignment
is arbitrary or whether it follows necessarily from the structural
properties of the integers 1, 2, and 3 themselves.

This paper argues the latter. We show that the cardinality-1,
cardinality-2, and cardinality-3 structures, examined across five
independent mathematical domains (Von Neumann ordinal theory, small
category theory, graph theory, simplex topology, and the theory of
symmetric groups), each exhibit a qualitative transition at the
corresponding integer that is isomorphic to the role transition
\(N \to D \to C\). We further show that this mapping is independently
recovered, without knowledge of the tholonic framework, by Peirce's
phenomenological categories (Firstness, Secondness, Thirdness), Kant's
categories of quantity (Unity, Plurality, Totality), and the distinction
calculus of Spencer-Brown. The convergence of five formal domains and
three independent philosophical traditions on the same mapping at the
same integers constitutes, we argue, strong evidence that the role
assignment is not arbitrary but follows from what the integers
structurally are.

We do not claim a formal theorem in the strongest sense: we cannot rule
out that a future algebraic framework reorders or augments the mapping.
What we claim is more modest and more defensible: within the class of
minimal recursive systems satisfying the conditions of Lemma 4.2 of
\href{https://github.com/baardev/truevalue/raw/main/docnav/Research/papers/3_minimal-recursive-triadic-framework/3_minimal-recursive-triadic-framework.pdf}{paper
3}, the mapping from cardinality to role is the unique assignment
consistent with the structural properties of 1, 2, and 3 as they appear
across the five domains examined. Resolving this into a strict theorem
is identified as an open problem.

\begin{center}\rule{0.5\linewidth}{0.5pt}\end{center}

\subsection{1. Introduction}\label{introduction}

\subsubsection{1.1 The gap in the prior
argument}\label{the-gap-in-the-prior-argument}

\href{https://github.com/baardev/truevalue/raw/main/docnav/Research/papers/3_minimal-recursive-triadic-framework/3_minimal-recursive-triadic-framework.pdf}{Paper
3 of this series} establishes that three functional roles, a negotiation
state, a defining/bounding auxiliary, and a contributing/integrating
auxiliary, are the minimum necessary for non-trivial convergent
recursion (Lemma 4.2).
\href{https://github.com/baardev/truevalue/raw/main/docnav/Research/papers/3_minimal-recursive-triadic-framework/3_minimal-recursive-triadic-framework.pdf}{Paper
3} also establishes that ``these roles are assigned by function, not by
numerical value''
\href{https://github.com/baardev/truevalue/raw/main/docnav/Research/papers/3_minimal-recursive-triadic-framework/3_minimal-recursive-triadic-framework.pdf}{Mil26b},
§4.2. This is an important caution: two variables may share the same
numerical seed while playing structurally distinct roles. The roles are
real; the numbers that happen to inhabit them at any moment are
incidental.

But there is a deeper question that
\href{https://github.com/baardev/truevalue/raw/main/docnav/Research/papers/3_minimal-recursive-triadic-framework/3_minimal-recursive-triadic-framework.pdf}{paper
3} does not address. The roles themselves are labeled by position: first
variable, second variable, third variable. When we write the triad as
\((N, D, C)\), we are placing \(N\) in position 1, \(D\) in position 2,
and \(C\) in position 3. Is this ordering arbitrary, a mere notational
convention, or does something about being-in-position-1 versus
being-in-position-3 actually determine the role that a variable must
play?

Put directly: is the negotiation role assigned to position 1
\emph{because} position 1 corresponds to the integer 1, and the integer
1 has inherent structural properties that are the same as what
negotiation functionally is? If so, the role assignment is not a
labeling convention but a consequence of the nature of the integers.
This paper argues that it is.

\subsubsection{1.2 Why this matters}\label{why-this-matters}

The practical consequences are significant for the entire framework. If
role assignment is merely conventional, then the tholonic framework is a
powerful organizational taxonomy: it groups recursive systems into a
three-role schema that happens to be useful. If role assignment is
necessary, the framework is something stronger: a consequence of the
inherent structure of the first three positive integers, and therefore
applicable without stipulation to any system that instantiates those
integers in their natural order.

The difference is the difference between a useful notation and a natural
law.

\subsubsection{1.3 Approach and scope}\label{approach-and-scope}

We proceed in three stages. First, we examine what the integers 1, 2,
and 3 structurally are across five independent mathematical domains,
showing that in each domain, the qualitative character of the structure
of cardinality \(n\) exhibits a transition at \(n = 1\), \(n = 2\), and
\(n = 3\) that mirrors the role transitions \(N\), \(D\), \(C\). Second,
we examine three independent philosophical traditions that arrive at the
same qualitative mapping without knowledge of the tholonic framework.
Third, we state the thesis precisely, characterize what would falsify
it, identify open problems, and connect the result to the broader
series.

\textbf{What this paper claims:} The mapping from cardinality to triadic
role is the unique assignment consistent with the structural properties
of 1, 2, and 3 in the five domains examined, and it is independently
confirmed by three philosophical traditions.

\textbf{What this paper does not claim:} That a formal theorem in the
sense of a derivation from Peano axioms alone establishes this; that no
other role assignments are conceivable in different logical systems; or
that the tholonic framework is the only framework in which this mapping
is relevant.

The tone throughout is argument and evidence, not proof by fiat. The
reader is invited to evaluate each domain independently and judge
whether the convergence is coincidental or structural.

\begin{center}\rule{0.5\linewidth}{0.5pt}\end{center}

\subsection{2. Historical and Philosophical
Context}\label{historical-and-philosophical-context}

Before the formal argument, we note that multiple intellectual
traditions have independently assigned qualitative meanings to the first
three positive integers. We present these not as authorities to be
deferred to, but as evidence that independent analysis consistently
recovers the same mapping, which would be surprising if the mapping were
merely conventional.

\subsubsection{2.1 Pythagorean monad, dyad, and
triad}\label{pythagorean-monad-dyad-and-triad}

The Pythagorean tradition, as reconstructed from Aristotle's
\emph{Metaphysics} (Book A, 985b--986a) and Nicomachus of Gerasa's
\emph{Introduction to Arithmetic} (c.~100 CE), assigned qualitative
character to the first integers as follows {[}Nic00{]}:

\begin{itemize}
\tightlist
\item
  The \textbf{monad} (1): the principle of unity; self-same, not yet
  divided; the generator from which all other numbers proceed but which
  is not itself generated. Nicomachus calls it ``the fountainhead and
  root of all numbers.'' It corresponds to identity without
  differentiation.
\item
  The \textbf{dyad} (2): the first departure from unity; the principle
  of opposition, otherness, and indefinite duality. Aristotle reports
  that Plato identified the dyad with the ``great and small,'' the
  principle of unlimited variance. It corresponds to the establishment
  of a limit or boundary between two things.
\item
  The \textbf{triad} (3): the first complete number in the sense of
  having a beginning, middle, and end; the first odd number that cannot
  be halved; the first number that enables a return. In Pythagorean
  arithmetic, three is the number of harmony: it combines the monad's
  unity and the dyad's distinction into a synthesis.
\end{itemize}

The correspondence to \((N, D, C)\) is direct: monad as negotiation
(unity, the running state), dyad as definition (the first limit, the
boundary), triad as contribution (the synthesis that enables return).

\subsubsection{2.2 Kant's categories of
quantity}\label{kants-categories-of-quantity}

Kant's \emph{Critique of Pure Reason} (1781, A80/B106) presents twelve
categories organized in four groups of three. The first group, the
categories of \textbf{Quantity}, is {[}Kan81{]}:

\begin{itemize}
\tightlist
\item
  \textbf{Unity} (corresponding to the integer 1 in Kant's own table)
\item
  \textbf{Plurality} (corresponding to 2)
\item
  \textbf{Totality} (corresponding to 3, which Kant explicitly describes
  as ``the synthesis of unity and plurality'')
\end{itemize}

Kant's derivation is transcendental: these are the a priori forms
through which any object can be cognized as to its extent. Unity is the
condition for treating something as a single thing. Plurality is the
condition for treating things as many. Totality is, critically, not
merely ``many unities'' but the synthesis: the condition for grasping a
manifold as a complete whole.

This is precisely the \((N, D, C)\) structure: Unity = \(N\) (the
running single state), Plurality = \(D\) (the bounding differentiation
that establishes what is more than one), Totality = \(C\) (the
integrating synthesis). Kant arrived at this table through analysis of
the logical forms of judgment, entirely independently of any dynamical
or recursive consideration.

\subsubsection{2.3 Peirce's phenomenological
categories}\label{peirces-phenomenological-categories}

Charles Sanders Peirce, in his \emph{Collected Papers} (1931, vol.~1,
§§300-353) and throughout his phenomenological writings, argued for
three universal categories of experience, which he termed
\textbf{Firstness}, \textbf{Secondness}, and \textbf{Thirdness}
{[}Pei31{]}:

\begin{itemize}
\tightlist
\item
  \textbf{Firstness}: the quality of something considered purely in
  itself, without reference to anything else. It is monadic: pure
  feeling, pure potentiality, the state of a thing as it simply is.
  Peirce gives the example of a color or a sensation considered as a
  pure quale, before comparison with any other sensation.
\item
  \textbf{Secondness}: the quality of something as it stands in relation
  to one other thing. It is dyadic: brute reaction, resistance, the
  sheer ``otherness'' of the world. It requires two terms and is the
  category of existence as opposition.
\item
  \textbf{Thirdness}: the quality of something that brings two other
  things into relation \emph{through itself}. It is triadic and
  irreducibly so: representation, mediation, law. Peirce showed formally
  that triadic relations cannot be decomposed into dyadic ones without
  remainder, while all tetradic and higher relations can be decomposed
  into triads.
\end{itemize}

Peirce's Firstness, Secondness, and Thirdness map directly and precisely
to \(N\), \(D\), and \(C\): - \(N\) (negotiation, the running state)
\(\leftrightarrow\) Firstness (the state as it is in itself) - \(D\)
(definition, the bounding parameter) \(\leftrightarrow\) Secondness (the
dyadic limit, the ``other'' that bounds) - \(C\) (contribution, the
integrating parameter) \(\leftrightarrow\) Thirdness (the mediation that
synthesizes and enables recursion)

Peirce derived these categories through meticulous phenomenological
analysis over decades, entirely independently of the tholonic framework.
His result, that three irreducible categories are both necessary and
sufficient for the description of any phenomenon, is a significant
independent confirmation of the tholonic structure.

Crucially, Peirce's formal argument for the irreducibility of Thirdness
{[}Pei31, vol.~1, §346{]} is structurally analogous to Lemma 4.2 of
\href{https://github.com/baardev/truevalue/raw/main/docnav/Research/papers/3_minimal-recursive-triadic-framework/3_minimal-recursive-triadic-framework.pdf}{paper
3}: a dyadic system (two elements in mutual relation) cannot mediate or
represent; mediation requires a third term that is not reducible to
either of the two it relates.

\subsubsection{2.4 Hegel's dialectical
triad}\label{hegels-dialectical-triad}

Hegel's dialectic (thesis, antithesis, synthesis) in the \emph{Science
of Logic} (1812) {[}Heg12{]} maps as: - \textbf{Thesis} (1): the initial
position, posited in itself. \(\leftrightarrow N\) - \textbf{Antithesis}
(2): the negation, the bounding limit that defines the thesis by
opposing it. \(\leftrightarrow D\) - \textbf{Synthesis} (3): the
sublation (\emph{Aufhebung}) that preserves both while transcending the
opposition. \(\leftrightarrow C\)

The synthesis in Hegel is not merely the average of thesis and
antithesis. It is a higher-order integration that contains both while
enabling the process to continue at a new level. This is precisely what
the \(C\) variable does in the tholonic recurrence: it integrates past
states in a way that enables the next recursion.

\subsubsection{2.5 Spencer-Brown's distinction
calculus}\label{spencer-browns-distinction-calculus}

George Spencer-Brown's \emph{Laws of Form} (1969) {[}SB69{]} begins with
a single primitive operation, the act of making a distinction. From this
primitive, two states arise: the marked state and the unmarked state.
The first law of form asserts that crossing the mark twice returns to
the origin. The formal system then has:

\begin{itemize}
\tightlist
\item
  \textbf{One mark}: the act of distinction itself; the state of being
  marked; pure assertion. \(\leftrightarrow N\) (1)
\item
  \textbf{Two states}: marked and unmarked; the first duality.
  \(\leftrightarrow D\) (2)
\item
  \textbf{Re-entry} (the system referring to itself through the
  distinction): the third element that makes the calculus
  self-referential and enables recursive depth. \(\leftrightarrow C\)
  (3)
\end{itemize}

Spencer-Brown's re-entry operator is the point at which the system first
becomes capable of self-reference. Without it, the calculus describes
only static distinction. With it, the calculus generates dynamics. This
is again the contribution of the third position: the enabling of
recursion. Spencer-Brown was explicit that re-entry is the minimum
augmentation required to move from a static to a dynamic logic.

\begin{center}\rule{0.5\linewidth}{0.5pt}\end{center}

\subsection{3. Structural Properties of 1, 2, and 3 in Five Mathematical
Domains}\label{structural-properties-of-1-2-and-3-in-five-mathematical-domains}

We now turn from philosophy to mathematics. In each of five domains, we
characterize what cardinality-1, cardinality-2, and cardinality-3
structures are, and show that the qualitative transitions at each
cardinality are isomorphic to the role transitions \(N \to D \to C\).

\subsubsection{3.1 Von Neumann ordinals}\label{von-neumann-ordinals}

The Von Neumann construction defines the natural numbers as sets:

\[0 = \emptyset, \quad 1 = \{0\} = \{\emptyset\}, \quad 2 = \{0,1\} = \{\emptyset, \{\emptyset\}\}, \quad 3 = \{0,1,2\}, \quad \ldots\]

Each integer \(n\) is the set of all its predecessors. This definition
is not merely a labeling convention; it reveals the intrinsic structural
character of each integer as a set.

\textbf{The integer 1} is \(\{\emptyset\}\): the set containing exactly
one element (the empty set, i.e., nothing). Its only member relation is
\(0 \in 1\). The set 1 has exactly one element, one subset relation, and
one element-of relation. It is internally consistent, self-contained,
and has no internal differentiation. It cannot be partitioned into two
non-empty subsets (since it has only one element). Structurally, 1 is a
unity: it asserts itself without differentiation. This is the
negotiation role: the state as it simply is, evaluated for internal
coherence alone.

\textbf{The integer 2} is \(\{0, 1\}\): the first set with two distinct
members. For the first time, an internal distinction exists:
\(0 \neq 1\). The set 2 supports a non-trivial binary relation
(\(0 < 1\)). It can be partitioned into two non-empty subsets (\(\{0\}\)
and \(\{1\}\)). The power set of 2 has four elements:
\(\emptyset, \{0\}, \{1\}, \{0,1\}\), which includes for the first time
a ``bounded'' subset (the strict subset \(\{0\}\) is bounded by the
complement \(\{1\}\)). The qualitative transition at 2 is the emergence
of internal boundary: something (0) is now defined against something
else (1). This is the definition/limitation role: the first
establishment of a boundary.

\textbf{The integer 3} is \(\{0, 1, 2\}\): the first set that supports a
non-degenerate ordering with a middle element. The element \(1\) is
strictly between \(0\) and \(2\): \(0 < 1 < 2\). More significantly, 3
is the smallest set that supports a transitive closure that is not
trivial: the relation \(0 < 1\) and \(1 < 2\) forces \(0 < 2\) by
transitivity, creating a composed relation not directly stated. This
composition is the key structural property of 3: it is the smallest
cardinality at which the \emph{integration of prior relations into new
relations} is forced. This is the contribution/integration role: the
accumulation of prior steps into a composed output.

\textbf{Summary (Von Neumann):} 1 asserts without differentiating. 2
differentiates without composing. 3 composes without needing higher
cardinality. The role mapping is structural, not conventional.

\subsubsection{3.2 Small category theory}\label{small-category-theory}

A category with \(n\) objects (and identity morphisms only) is the
discrete category \(\mathbf{n}\). But more instructively, the ordinal
categories \(\mathbf{1}\), \(\mathbf{2}\), \(\mathbf{3}\) are defined as
follows {[}ML98{]}:

\begin{itemize}
\item
  \(\mathbf{1}\): one object, one morphism (the identity). No
  non-trivial arrows. No composition needed or possible. The terminal
  category: every category has exactly one functor to \(\mathbf{1}\). It
  is the categorical representation of unity: the unique ``view from
  nowhere'' that everything maps to.
\item
  \(\mathbf{2}\): two objects \(A, B\) and one non-trivial arrow
  \(f: A \to B\) (plus identities). The simplest non-trivial category.
  It represents the category of a single directed distinction: a
  relation from something to something-else. Every morphism in any
  category ``factors through'' a version of \(\mathbf{2}\) (each
  individual arrow is a functor from \(\mathbf{2}\)). It is the
  categorical representation of a boundary or limit: one thing in
  relation to exactly one other.
\item
  \(\mathbf{3}\): three objects \(A, B, C\) and arrows \(f: A \to B\),
  \(g: B \to C\), and their composite \(g \circ f: A \to C\) (plus
  identities). This is the first category in which non-trivial
  \textbf{composition} is required. The arrow \(g \circ f\) is not
  explicitly listed as a primitive but must exist by the axioms of
  category theory. The category \(\mathbf{3}\) is thus the first
  category that generates something new from its primitive data: the
  composite \(g \circ f\) is the integrating contribution of the two
  prior arrows. Functors from \(\mathbf{3}\) into other categories are
  exactly the composable pairs of morphisms.
\end{itemize}

In category theory, \(\mathbf{1}\) is the object, \(\mathbf{2}\) is the
morphism (the directed distinction), and \(\mathbf{3}\) is the
composable pair, the minimum structure in which something genuinely new
(the composite) is generated from two prior things. This is the \(N\),
\(D\), \(C\) structure exactly.

\textbf{Remark on higher cardinalities.} Categories
\(\mathbf{4}, \mathbf{5}, \ldots\) introduce additional composition data
but no new qualitative transitions: they add more composable sequences
without changing the structural character of composition itself. The
qualitative jump is complete at \(\mathbf{3}\).

\subsubsection{3.3 Graph theory}\label{graph-theory}

The complete graphs \(K_1\), \(K_2\), \(K_3\) exhibit the following
structural properties:

\begin{itemize}
\item
  \(K_1\): one vertex, no edges. A single node with no connections. It
  has no non-trivial subgraph. It represents a self-contained unit: no
  relation to any other. Structurally, it is the ``state in isolation,''
  the unit evaluated without reference to an external.
\item
  \(K_2\): two vertices, one edge. The first graph with a non-trivial
  relation. It supports exactly one path of length 1 and no cycles. It
  represents a binary relation (this connected to that) without any
  closure. It is the minimum graph for a directed distinction.
\item
  \(K_3\): three vertices, three edges. The \textbf{minimum graph that
  contains a cycle}: the path \(A \to B \to C \to A\). This is the
  qualitative transition: \(K_3\) is the smallest complete graph with
  non-zero cyclomatic complexity (\(\mu = E - V + 1 = 3 - 3 + 1 = 1\)).
  A cycle is a closed path: a sequence that \emph{returns to its
  origin}. Return is the enabling condition for recursion. Without a
  cycle, a graph can only propagate forward; it cannot feed back,
  integrate, or self-refer.
\end{itemize}

The contribution of \(K_3\)'s third vertex is specifically to close the
open path of \(K_2\) into a cycle. Removing any vertex from \(K_3\)
eliminates the cycle entirely (leaving \(K_2\)). The third vertex is
therefore the structural contributor that makes closure possible.

\textbf{Cyclomatic complexity progression:}
\[\mu(K_1) = 0, \quad \mu(K_2) = 0, \quad \mu(K_3) = 1, \quad \mu(K_4) = 4, \ldots\]

The first nonzero cyclomatic complexity appears exactly at \(K_3\). This
is not a coincidence but a consequence of the Euler formula for
connected graphs: a tree on \(n\) vertices has \(n-1\) edges and
\(\mu = 0\); the first additional edge creates the first cycle and
requires at least \(n = 3\) vertices to be non-degenerate.

\subsubsection{3.4 Simplex topology}\label{simplex-topology}

The \(k\)-simplex is the generalization of a triangle to \(k\)
dimensions {[}Hat02{]}:

\begin{itemize}
\item
  The \textbf{0-simplex} (\(k=0\)): a single point. No interior, no
  edges, no faces. Contractible to a point. Topologically, it is the
  simplest possible object: self-contained, without boundary structure.
\item
  The \textbf{1-simplex} (\(k=1\)): a line segment with two endpoints.
  It has a boundary (the two endpoints) but no interior faces. It
  represents a directed distinction with two endpoints. The boundary
  operator applied to a 1-simplex yields its two endpoint 0-simplices.
\item
  The \textbf{2-simplex} (\(k=2\)): a triangle with three vertices,
  three edges, and one 2-face (interior). It is the smallest simplex
  with an interior: the first topological object with non-trivial area.
  The \textbf{Euler characteristic} of the 2-simplex is
  \(\chi = V - E + F = 3 - 3 + 1 = 1\). More significantly, the boundary
  \(\partial\) of the 2-simplex (the triangular loop) is the first
  non-trivial boundary that is itself a closed cycle:
  \(\partial(\Delta^2) = [v_1, v_2] - [v_0, v_2] + [v_0, v_1]\), a sum
  of three 1-simplices forming a closed loop.
\end{itemize}

\href{https://github.com/baardev/truevalue/raw/main/docnav/Research/papers/3_minimal-recursive-triadic-framework/3_minimal-recursive-triadic-framework.pdf}{Paper
3 of this series} already uses the 2-simplex as the geometric substrate
for the triadic role partition
\href{https://github.com/baardev/truevalue/raw/main/docnav/Research/papers/3_minimal-recursive-triadic-framework/3_minimal-recursive-triadic-framework.pdf}{Mil26b},
§3. What the present paper adds is the observation that the progression
\(\Delta^0 \to \Delta^1 \to \Delta^2\) exhibits exactly the qualitative
transitions of unity, boundary, and closed integration: the 0-simplex is
a unit, the 1-simplex establishes a boundary (two vertices, one edge, no
closure), and the 2-simplex achieves closure (three vertices, three
edges, one interior face, one closed boundary cycle).

The 2-simplex is the minimum simplex that is not contractible to a
lower-dimensional object while preserving its full topological
character. In this sense, it is the minimum structure that
``contributes'' a genuinely new dimension.

\subsubsection{3.5 Symmetric groups}\label{symmetric-groups}

The symmetric group \(S_n\) is the group of all permutations of \(n\)
elements. The structural properties of \(S_1\), \(S_2\), \(S_3\) exhibit
clear qualitative transitions {[}Arm88{]}:

\begin{itemize}
\item
  \(S_1\): the trivial group, \(|S_1| = 1\), with only the identity
  permutation. No non-trivial structure whatsoever. It cannot represent
  any interesting symmetry because there is only one way to permute a
  single element.
\item
  \(S_2\): the cyclic group \(\mathbb{Z}_2\) of order 2, \(|S_2| = 2\).
  Two elements: the identity and the transposition (swap). This is the
  simplest non-trivial group: it has one non-trivial operation, and that
  operation is its own inverse. It represents the minimal symmetry of
  exchange, the binary flip. It is Abelian (commutative): the order of
  operations does not matter, because there is only one non-trivial
  operation.
\item
  \(S_3\): the first non-Abelian group, \(|S_3| = 6\). It is the
  smallest group in which \(ab \neq ba\) for some elements \(a, b\).
  Non-commutativity is a qualitative threshold: it means that the order
  of operations matters, which is the first sign of genuine complexity
  in the algebraic sense. \(S_3\) also contains subgroups of two
  different orders (1, 2, and 3), making it the first group with
  non-trivial subgroup structure of more than one type.
\end{itemize}

The qualitative progression is: \(S_1\) (no structure), \(S_2\)
(structure but trivially simple and commutative), \(S_3\) (first
genuinely complex structure with non-Abelian behavior). Non-Abelian
structure means that elements interact in ways that depend on their
order of application, which is the algebraic counterpart of integration:
the result depends on the history of how contributions have been
combined.

\subsubsection{3.6 Summary of the five
domains}\label{summary-of-the-five-domains}

\begin{longtable}[]{@{}
  >{\raggedright\arraybackslash}p{(\linewidth - 6\tabcolsep) * \real{0.0930}}
  >{\raggedright\arraybackslash}p{(\linewidth - 6\tabcolsep) * \real{0.3023}}
  >{\raggedright\arraybackslash}p{(\linewidth - 6\tabcolsep) * \real{0.2907}}
  >{\raggedright\arraybackslash}p{(\linewidth - 6\tabcolsep) * \real{0.3140}}@{}}
\toprule\noalign{}
\begin{minipage}[b]{\linewidth}\raggedright
Domain
\end{minipage} & \begin{minipage}[b]{\linewidth}\raggedright
\(n = 1\) (N: negotiation)
\end{minipage} & \begin{minipage}[b]{\linewidth}\raggedright
\(n = 2\) (D: definition)
\end{minipage} & \begin{minipage}[b]{\linewidth}\raggedright
\(n = 3\) (C: contribution)
\end{minipage} \\
\midrule\noalign{}
\endhead
\bottomrule\noalign{}
\endlastfoot
Von Neumann ordinals & Unity: no internal distinction & First internal
boundary & First transitive composition \\
Category theory & Terminal object: maps to everything & Directed
morphism: first non-trivial arrow & First forced composition: something
new generated \\
Graph theory & No edges, no relations & One edge, one relation, no cycle
& First cycle: recursion enabled \\
Simplex topology & Point: no boundary & Segment: boundary but no closure
& Triangle: first closed interior \\
Symmetric groups & Trivial group & \(\mathbb{Z}_2\): first non-trivial,
commutative & \(S_3\): first non-Abelian, first genuine complexity \\
\end{longtable}

The five domains are independent. Each uses different mathematical
objects and different definitions of structure. Yet in all five, the
qualitative character of cardinality 1 is unity/self-containment, the
qualitative character of cardinality 2 is the establishment of a
directed boundary or distinction, and the qualitative character of
cardinality 3 is the first closure, composition, or integration that
generates something beyond what the prior two steps contained.

This convergence is not a theorem in the narrow sense: we have not
derived it from a single set of axioms. It is, rather, an observation
across independent domains that the qualitative transitions in each
domain occur at the same integers and exhibit the same qualitative
character. We now argue that this convergence is evidence for the
thesis.

\begin{center}\rule{0.5\linewidth}{0.5pt}\end{center}

\subsection{4. The Core Argument}\label{the-core-argument}

\subsubsection{4.1 Statement}\label{statement}

\textbf{Thesis.} Let \(\mathcal{T}\) be any minimal recursive system
satisfying the conditions of Lemma 4.2 of
\href{https://github.com/baardev/truevalue/raw/main/docnav/Research/papers/3_minimal-recursive-triadic-framework/3_minimal-recursive-triadic-framework.pdf}{paper
3} (a distinguished state variable and two functionally independent
auxiliary variables). Assign to each variable a cardinality: 1 to the
distinguished state variable, 2 to the bounding auxiliary, 3 to the
integrating auxiliary. Then the functional role of each variable in
\(\mathcal{T}\) is isomorphic to the structural character of the
corresponding cardinality as described in §3.

Equivalently: the negotiation role is not assigned to position 1 by
convention. Position 1 \emph{is} structurally what the negotiation role
functionally requires (unity, self-containment, evaluation without
external reference). Position 2 \emph{is} structurally what the
definition role functionally requires (first directed distinction, the
establishment of a boundary). Position 3 \emph{is} structurally what the
contribution role functionally requires (first closure, composition,
integration of prior elements into something new).

\subsubsection{4.2 Argument by structural
isomorphism}\label{argument-by-structural-isomorphism}

The argument proceeds by structural isomorphism across the five domains
of §3 and the role definitions of
\href{https://github.com/baardev/truevalue/raw/main/docnav/Research/papers/3_minimal-recursive-triadic-framework/3_minimal-recursive-triadic-framework.pdf}{paper
3}.

\textbf{For the negotiation role \(N\) and the integer 1:}

The negotiation variable is ``the running state being iteratively
refined; the emergent quantity''
\href{https://github.com/baardev/truevalue/raw/main/docnav/Research/papers/3_minimal-recursive-triadic-framework/3_minimal-recursive-triadic-framework.pdf}{Mil26b},
Definition 4.3. It is evaluated for internal coherence: given the
current state, is the claim consistent with itself? It does not refer to
an external bound or an external corroboration. It simply asserts its
current value and is assessed on internal grounds.

This is precisely the structural character of cardinality 1 across all
five domains: \(\{\emptyset\}\) (no internal distinction),
\(\mathbf{1}\) (terminal category, no non-trivial arrows), \(K_1\) (no
edges), \(\Delta^0\) (no boundary), \(S_1\) (only the identity). In
every case, the cardinality-1 structure is self-contained,
self-referential in a trivial sense (it refers only to itself), and has
no internal differentiation. The negotiation role requires the same: it
evaluates a single position without reference to an external limit.

\textbf{For the definition role \(D\) and the integer 2:}

The definition variable is ``the bounding or constraining parameter;
what limits the state at each step''
\href{https://github.com/baardev/truevalue/raw/main/docnav/Research/papers/3_minimal-recursive-triadic-framework/3_minimal-recursive-triadic-framework.pdf}{Mil26b},
Definition 4.3. It establishes what the state is \emph{against}: the
audit scope, the policy boundary, the counterparty distinction. It is
the first externality.

This is precisely the structural character of cardinality 2 across all
five domains: \(\{0,1\}\) (first internal distinction), \(\mathbf{2}\)
(one non-trivial directed arrow, the first morphism), \(K_2\) (one edge,
one binary relation), \(\Delta^1\) (two endpoints, a directed boundary),
\(S_2\) (one non-trivial operation, the transposition). In every case,
the cardinality-2 structure introduces the first external reference:
something in relation to one other thing. The definition role requires
the same: it bounds the state by referencing something outside the state
itself.

\textbf{For the contribution role \(C\) and the integer 3:}

The contribution variable is ``the accumulating or synthesizing
parameter; what drives growth or combines past states''
\href{https://github.com/baardev/truevalue/raw/main/docnav/Research/papers/3_minimal-recursive-triadic-framework/3_minimal-recursive-triadic-framework.pdf}{Mil26b},
Definition 4.3. It integrates: it takes what has happened (the prior
state, the prior boundary) and produces something new from the
combination. It is the role that makes recursion possible, because
recursion requires that the output of one step feed into the next.

This is precisely the structural character of cardinality 3 across all
five domains: first transitive composition (Von Neumann), first forced
composite morphism (\(\mathbf{3}\)), first cycle (\(K_3\), cyclomatic
complexity = 1), first closed interior (\(\Delta^2\)), first non-Abelian
group (\(S_3\)). In every case, cardinality 3 is the first cardinality
at which something genuinely new is generated from the combination of
prior elements: a composed path, a forced morphism, a closed loop, an
interior, a non-trivial subgroup structure. The contribution role
requires the same: it generates a new state from the combination of
\(N\) and \(D\).

\subsubsection{4.3 Uniqueness of the
mapping}\label{uniqueness-of-the-mapping}

A natural question is whether the mapping could be permuted: could \(D\)
be at position 1 and \(N\) at position 2, for example? The structural
argument says no, for the following reason.

The definition role requires an external reference: something to be
defined against. A cardinality-1 structure has no external reference by
definition (it has one element, no internal distinction, no other).
Therefore the definition role cannot inhabit position 1. It requires at
least cardinality 2.

The contribution role requires a composition: it must combine two prior
elements into a new one. A cardinality-2 structure has two elements but
no composition (no third element into which the two combine). Therefore
the contribution role cannot inhabit position 2. It requires at least
cardinality 3.

The negotiation role requires only internal coherence: no external
reference, no composition. A cardinality-1 structure satisfies this
exactly. Therefore the negotiation role inhabits position 1.

The mapping \((N \to 1, D \to 2, C \to 3)\) is thus the unique
assignment consistent with the structural requirements of each role and
the structural properties of each cardinality. Any permutation violates
at least one structural constraint.

\begin{center}\rule{0.5\linewidth}{0.5pt}\end{center}

\subsection{5. Connection to the Tholonic
Series}\label{connection-to-the-tholonic-series}

\subsubsection{\texorpdfstring{5.1
\href{https://github.com/baardev/truevalue/raw/main/docnav/Research/papers/1_recursive-tholonic-five-constants/1_recursive-tholonic-five-constants.pdf}{Paper
1} (five
constants)}{5.1 Paper 1 (five constants)}}\label{paper-1-five-constants}

The five branches of the tholonic ladder recurrence
\href{https://github.com/baardev/truevalue/raw/main/docnav/Research/papers/1_recursive-tholonic-five-constants/1_recursive-tholonic-five-constants.pdf}{Mil26a}
use \((N, D, C)\) triples with the role structure established here. The
\(\pi/4\) branch uniquely requires seeds \((1, 3, 5)\), where
\(N_0 = 1\) (the multiplicative identity, the unit), \(D_0 = 3\) and
\(C_0 = 5\) (the first and second odd integers beyond 1, corresponding
to the Definition and Contribution axis multipliers of the tholonic
geometry). The assignment of \(N_0 = 1\) to the negotiation position is
thus not arbitrary: the negotiation variable begins at 1 because 1 is
what the negotiation role structurally is.

\subsubsection{\texorpdfstring{5.2
\href{https://github.com/baardev/truevalue/raw/main/docnav/Research/papers/2_supply-chain-transparency-tvpci/2_supply-chain-transparency-tvpci.pdf}{Paper
2} (TVPCI supply-chain
transparency)}{5.2 Paper 2 (TVPCI supply-chain transparency)}}\label{paper-2-tvpci-supply-chain-transparency}

The supply-chain framework assigns \(N_p\) (the custody state), \(D_p\)
(the bounding evidence), and \(C_p\) (the corroboration evidence) to
each phase
\href{https://github.com/baardev/truevalue/raw/main/docnav/Research/papers/2_supply-chain-transparency-tvpci/2_supply-chain-transparency-tvpci.pdf}{Mil26c}.
Section 4.1 of that paper now states that ``the integers carry inherent
qualitative meaning rooted in their structural properties.'' The present
paper is the formal grounding for that claim.

\subsubsection{5.3 Papers 3 and 4}\label{papers-3-and-4}

\href{https://github.com/baardev/truevalue/raw/main/docnav/Research/papers/3_minimal-recursive-triadic-framework/3_minimal-recursive-triadic-framework.pdf}{Paper
3} establishes that three roles are necessary. The present paper
establishes why those three roles correspond to cardinalities 1, 2, 3.
\href{https://github.com/baardev/truevalue/raw/main/docnav/Research/papers/4_game-theoretic-triadic-balance/4_game-theoretic-triadic-balance.pdf}{Paper
4}
\href{https://github.com/baardev/truevalue/raw/main/docnav/Research/papers/4_game-theoretic-triadic-balance/4_game-theoretic-triadic-balance.pdf}{Mil26d}
recasts \(D\) and \(C\) as opposing strategic agents: \(D\) contracting
(bounding) and \(C\) expanding (contributing). The game-theoretic
formulation is consistent with the structural argument here: the
bounding player requires two elements (a move and a counter-move, the
cardinality-2 relation), and the integrating player requires three
elements (a prior state, a prior bound, and the synthesis).

\subsubsection{\texorpdfstring{5.4
\href{https://github.com/baardev/truevalue/raw/main/docnav/Research/papers/5_tholonic-twistor-connection/5_tholonic-twistor-connection.pdf}{Paper
5} (twistor
connection)}{5.4 Paper 5 (twistor connection)}}\label{paper-5-twistor-connection}

The connection to twistor geometry
\href{https://github.com/baardev/truevalue/raw/main/docnav/Research/papers/5_tholonic-twistor-connection/5_tholonic-twistor-connection.pdf}{Mil26e}
raises the possibility that the \((N, D, C)\) structure maps onto
fundamental objects in theoretical physics. If the twistor
correspondence is genuine, then the qualitative nature of 1, 2, 3 as
role-bearing integers would be a structural property at the level of
mathematical physics, not merely of recursive taxonomy.

\begin{center}\rule{0.5\linewidth}{0.5pt}\end{center}

\subsection{6. Independent Confirmation: Convergence of
Traditions}\label{independent-confirmation-convergence-of-traditions}

We now observe that the five mathematical domains of §3 and the four
philosophical traditions of §2 are independent in the following precise
sense: none of the authors cited in §2 had access to the tholonic
framework, none of the mathematical results of §3 were derived to
support the tholonic framework, and the tholonic framework was not
derived by surveying these precedents. The convergence is therefore
evidential: if the mapping were merely conventional, there is no reason
to expect it to be recovered by Kant's transcendental analytic, Peirce's
phenomenology, Spencer-Brown's distinction calculus, Pythagorean
arithmetic, and five branches of mathematics independently.

The convergence is summarized:

\begin{longtable}[]{@{}
  >{\raggedright\arraybackslash}p{(\linewidth - 6\tabcolsep) * \real{0.4706}}
  >{\raggedright\arraybackslash}p{(\linewidth - 6\tabcolsep) * \real{0.1765}}
  >{\raggedright\arraybackslash}p{(\linewidth - 6\tabcolsep) * \real{0.1765}}
  >{\raggedright\arraybackslash}p{(\linewidth - 6\tabcolsep) * \real{0.1765}}@{}}
\toprule\noalign{}
\begin{minipage}[b]{\linewidth}\raggedright
Source
\end{minipage} & \begin{minipage}[b]{\linewidth}\raggedright
1
\end{minipage} & \begin{minipage}[b]{\linewidth}\raggedright
2
\end{minipage} & \begin{minipage}[b]{\linewidth}\raggedright
3
\end{minipage} \\
\midrule\noalign{}
\endhead
\bottomrule\noalign{}
\endlastfoot
Pythagorean tradition & Monad (unity, self-same) & Dyad (opposition,
limit) & Triad (harmony, synthesis) \\
Kant (Critique) & Unity & Plurality & Totality (synthesis of unity and
plurality) \\
Peirce (categories) & Firstness (monadic, pure quality) & Secondness
(dyadic, reaction, otherness) & Thirdness (triadic, mediation, law) \\
Spencer-Brown & The mark (pure assertion) & Marked/unmarked (first
duality) & Re-entry (self-reference, recursion) \\
Hegel & Thesis (position in itself) & Antithesis (negation, limit) &
Synthesis (sublation, integration) \\
Von Neumann ordinals & Unity: no internal distinction & First boundary &
First composition \\
Category theory & Terminal object & First morphism & First forced
composition \\
Graph theory & No edges & One edge, no cycle & First cycle \\
Simplex topology & Point & Segment & Triangle (first interior, first
closure) \\
Symmetric groups & Trivial group & \(\mathbb{Z}_2\) (commutative) &
\(S_3\) (non-Abelian, first complexity) \\
\textbf{Tholonic framework} & \(N\): negotiation, unity state & \(D\):
definition, bounding & \(C\): contribution, integration \\
\end{longtable}

Eleven independent sources; one mapping.

\begin{center}\rule{0.5\linewidth}{0.5pt}\end{center}

\subsection{7. Objections and Responses}\label{objections-and-responses}

\subsubsection{7.1 ``The mapping is post-hoc
rationalization''}\label{the-mapping-is-post-hoc-rationalization}

\emph{Objection.} You began with \(N\), \(D\), \(C\) and then searched
for mathematical structures that match. Any three distinct roles could
be matched post-hoc to any three mathematical structures by selecting
the right framing.

\emph{Response.} The objection would be compelling if the mathematical
structures had been selected to match. But the structures in §3 are not
selected: Von Neumann ordinals, ordinal categories, complete graphs,
simplices, and symmetric groups are among the most fundamental and
well-studied objects in mathematics, and their structural properties at
cardinalities 1, 2, and 3 are established results that precede the
tholonic framework by decades to centuries. The argument is not ``these
structures match \(N\), \(D\), \(C\)''; it is ``these structures all
exhibit the same qualitative transition pattern, and that pattern
matches \(N\), \(D\), \(C\), which were derived from an independent
recursive analysis.'' The convergence cannot be dismissed as selection
bias when the mathematical results are prior, independent, and standard.

Furthermore, the philosophical sources in §2 (Peirce, Kant,
Spencer-Brown) all arrived at the mapping through methods entirely
unrelated to the tholonic framework and to each other. Post-hoc
selection does not explain their independent convergence.

\subsubsection{7.2 ``The transitions at 1, 2, 3 are trivial because
these are the first three
integers''}\label{the-transitions-at-1-2-3-are-trivial-because-these-are-the-first-three-integers}

\emph{Objection.} Of course structures at cardinality 1, 2, 3 are
qualitatively different; any three successive integers would show
structural differences. There is nothing special about 1, 2, 3
specifically.

\emph{Response.} This objection fails for two reasons. First, the
qualitative transitions in §3 are not merely ``different at each step.''
They are specifically the transitions from self-containment to external
relation to composition: a three-step sequence that exactly matches the
three-role requirement of Lemma 4.2. The transitions at \(n = 4, 5, 6\)
do not produce new role classes of the same kind; they produce more
complex instances of cardinality-3 structure (more cycles, more
compositions, more non-Abelian complexity), not new roles.

Second, Peirce explicitly argued that Thirdness is the \emph{final}
irreducible category: all higher-arity relations can be decomposed into
triadic ones, but triadic relations cannot be decomposed into dyadic
ones {[}Pei31, vol.~1, §346{]}. This is a formal result, not a
stipulation. The qualitative sequence terminates at 3 in a
mathematically precise sense.

\subsubsection{7.3 ``Different axiom systems might produce a different
mapping''}\label{different-axiom-systems-might-produce-a-different-mapping}

\emph{Objection.} The structural properties described in §3 depend on
specific axioms (ZFC set theory, category theory as standardly
formulated, etc.). In a different foundational system, the mapping might
differ.

\emph{Response.} This is a fair caution, and we note it as a genuine
limitation (see §8). However, the convergence of five independent
mathematical domains under different axiom systems, all yielding the
same qualitative mapping, significantly reduces the force of this
objection. If the mapping were an artifact of a specific axiom system,
we would not expect it to appear in the category-theoretic,
graph-theoretic, and topological settings simultaneously, as these have
different foundational assumptions. The robustness of the mapping across
frameworks is itself evidence of its non-artificiality.

\subsubsection{7.4 ``Assigning quality to numbers is
numerology''}\label{assigning-quality-to-numbers-is-numerology}

\emph{Objection.} The claim that integers have inherent qualitative
meanings sounds like numerology or mysticism.

\emph{Response.} The distinction between numerology and the present
argument is precision and falsifiability. Numerology assigns meanings to
numbers by assertion, cultural convention, or symbolic association, and
the resulting claims are typically unfalsifiable. The present argument
assigns structural properties to integers by formal mathematical
analysis, and those assignments are falsifiable: if a formal argument
showed that a cardinality-1 structure can be made to exhibit the
bounding role, or that cardinality-2 structure can exhibit the
integrating role, the thesis would be challenged.

Furthermore, Peirce, Kant, and Spencer-Brown are not numerologists; they
are among the most rigorous analytical thinkers in the Western
tradition. Their independent convergence on the same mapping, arrived at
through logical and phenomenological analysis, is not numerological.

The thesis is not that ``the number 3 is holy'' or ``3 means
synthesis.'' The thesis is the bounded and falsifiable claim: within the
class of minimal recursive systems satisfying the conditions of Lemma
4.2, the unique role assignment consistent with the structural
properties of cardinalities 1, 2, and 3 is \((N, D, C)\).

\begin{center}\rule{0.5\linewidth}{0.5pt}\end{center}

\subsection{8. Open Problems}\label{open-problems}

The present paper argues the thesis through structural correspondence
and cross-domain convergence. Several formal gaps remain.

\textbf{Open problem 1 (Formal theorem from axioms).} Derive the
role-to-cardinality mapping as a theorem from the axioms of Lemma 4.2 of
\href{https://github.com/baardev/truevalue/raw/main/docnav/Research/papers/3_minimal-recursive-triadic-framework/3_minimal-recursive-triadic-framework.pdf}{paper
3} and the Peano axioms (or ZFC), without appeal to the cross-domain
convergence argument. This would elevate the thesis from a
well-evidenced claim to a formal result.

\textbf{Open problem 2 (Higher cardinalities).} Characterize the
structural properties of cardinalities \(n \ge 4\) in the context of
recursive systems. Does cardinality 4 introduce a fourth role, or does
it produce a higher-order instance of one of the existing three?
Peirce's result that Thirdness is the final irreducible category
suggests no new role appears at \(n = 4\), but a formal confirmation in
the tholonic setting would be valuable.

\textbf{Open problem 3 (The zero case).} The integer 0, as \(\emptyset\)
in Von Neumann ordinal theory, is the empty set: the state of absolute
non-distinction, prior even to unity. Does it correspond to a zeroth
role in the tholonic framework (the ground state, the unmanifest
potential), and if so, what are the consequences for the recurrence
structure?

\textbf{Open problem 4 (Cross-foundational robustness).} Verify the
mapping in foundational systems other than ZFC: homotopy type theory
(HoTT), constructive type theory, and paraconsistent logic. If the
mapping survives in systems with fundamentally different axioms, the
evidence for its non-artificiality strengthens further.

\textbf{Open problem 5 (Connection to physical constants).} If the five
classical constants (\(\pi/4\), \(\phi\), \(e\), \(\sqrt{2}\),
\(\ln 2\)) emerge from the \((N, D, C)\) structure, and the
\((N, D, C)\) structure follows from the qualitative properties of 1, 2,
3, then the constants themselves may be consequences of the qualitative
structure of the first three integers. Making this chain explicit would
be a significant result.

\begin{center}\rule{0.5\linewidth}{0.5pt}\end{center}

\subsection{9. Conclusion}\label{conclusion}

We have argued that the assignment of roles to positions in the tholonic
triad is not a notational convention. The negotiation role (\(N\),
position 1) corresponds to the integer 1 because cardinality-1
structures are, across five independent mathematical domains, precisely
what the negotiation role functionally requires: self-contained,
internally consistent, without external reference. The definition role
(\(D\), position 2) corresponds to the integer 2 because cardinality-2
structures are the first to exhibit a directed external relation, a
boundary, precisely what the definition role functionally requires. The
contribution role (\(C\), position 3) corresponds to the integer 3
because cardinality-3 structures are the first to exhibit composition,
closure, and the generation of something new from prior elements,
precisely what the contribution role functionally requires.

This structural isomorphism is independently confirmed by eleven sources
across two methodological traditions (formal mathematics and
philosophical analysis), none of which was constructed to support the
tholonic framework.

The practical consequence is that the tholonic framework is not merely a
useful organizational schema. It is an expression of the structural
properties of the first three positive integers as they manifest in any
system complex enough to exhibit recursion. A system with three roles in
the \((N, D, C)\) pattern is not using a clever taxonomy; it is
exhibiting the minimum structure that the integers themselves make
available for non-trivial self-organization.

This does not make the framework immune to revision. Open problems 1
through 5 identify genuine formal gaps. But it does mean that the
framework's foundation is not arbitrary, and that the convergence of
mathematics, phenomenology, and the tholonic series on the same mapping
at the same integers is not a coincidence to be dismissed but a
structural fact to be explained.

\begin{center}\rule{0.5\linewidth}{0.5pt}\end{center}

\subsection{References}\label{references}

\begin{enumerate}
\def\labelenumi{\arabic{enumi}.}
\tightlist
\item
  Armstrong, M. A. \emph{Groups and Symmetry.} Springer, 1988.
  (Symmetric groups \(S_n\) and their structural properties.)
\item
  Aristotle. \emph{Metaphysics.} Translated by W. D. Ross. Oxford
  University Press, 1924. (Book A, 985b--986a: Pythagorean number
  doctrine.)
\item
  Hatcher, A. \emph{Algebraic Topology.} Cambridge University Press,
  2002. (Simplicial complexes and Euler characteristic.)
\item
  Hegel, G. W. F. \emph{Science of Logic.} Translated by A. V. Miller.
  Allen and Unwin, 1969 {[}1812{]}. (Dialectical triad: thesis,
  antithesis, synthesis.)
\item
  Kant, I. \emph{Critique of Pure Reason.} Translated by N. K. Smith.
  Macmillan, 1929 {[}1781{]}. (A80/B106: Table of categories; categories
  of quantity: unity, plurality, totality.)
\item
  Mac Lane, S. \emph{Categories for the Working Mathematician.} 2nd
  ed.~Springer, 1998. (Ordinal categories \(\mathbf{1}\),
  \(\mathbf{2}\), \(\mathbf{3}\) and functors between them.)
\item
  Milton, J. W. \emph{Emergence of Classical Constants from a Minimal
  Recursive Triadic Framework.} Companion manuscript
  (\href{https://github.com/baardev/truevalue/raw/main/docnav/Research/papers/1_recursive-tholonic-five-constants/1_recursive-tholonic-five-constants.pdf}{paper
  1}), this repository, 2026.
  \href{https://github.com/baardev/truevalue/raw/main/docnav/Research/papers/1_recursive-tholonic-five-constants/1_recursive-tholonic-five-constants.pdf}{Mil26a}
\item
  Milton, J. W. \emph{Phase-Resolved Transparency Classification in
  Commodity Supply Chains: A Structural Triadic Scoring Framework
  (TVPCI).} Companion manuscript
  (\href{https://github.com/baardev/truevalue/raw/main/docnav/Research/papers/2_supply-chain-transparency-tvpci/2_supply-chain-transparency-tvpci.pdf}{paper
  2}), this repository, 2026.
  \href{https://github.com/baardev/truevalue/raw/main/docnav/Research/papers/2_supply-chain-transparency-tvpci/2_supply-chain-transparency-tvpci.pdf}{Mil26c}
\item
  Milton, J. W. \emph{A Minimal Recursive Triadic Framework for
  Self-Similar Hierarchical Systems.} Companion manuscript
  (\href{https://github.com/baardev/truevalue/raw/main/docnav/Research/papers/3_minimal-recursive-triadic-framework/3_minimal-recursive-triadic-framework.pdf}{paper
  3}), this repository, 2026.
  \href{https://github.com/baardev/truevalue/raw/main/docnav/Research/papers/3_minimal-recursive-triadic-framework/3_minimal-recursive-triadic-framework.pdf}{Mil26b}
\item
  Milton, J. W. \emph{Game-Theoretic Framing of the Triadic Balance
  Condition.} Companion manuscript
  (\href{https://github.com/baardev/truevalue/raw/main/docnav/Research/papers/4_game-theoretic-triadic-balance/4_game-theoretic-triadic-balance.pdf}{paper
  4}), this repository, 2026.
  \href{https://github.com/baardev/truevalue/raw/main/docnav/Research/papers/4_game-theoretic-triadic-balance/4_game-theoretic-triadic-balance.pdf}{Mil26d}
\item
  Milton, J. W. \emph{Tholonic-Twistor Connection.} Companion manuscript
  (\href{https://github.com/baardev/truevalue/raw/main/docnav/Research/papers/5_tholonic-twistor-connection/5_tholonic-twistor-connection.pdf}{paper
  5}), this repository, 2026.
  \href{https://github.com/baardev/truevalue/raw/main/docnav/Research/papers/5_tholonic-twistor-connection/5_tholonic-twistor-connection.pdf}{Mil26e}
\item
  Nicomachus of Gerasa. \emph{Introduction to Arithmetic.} Translated by
  M. L. D'Ooge. Macmillan, 1926 {[}c.~100 CE{]}. (Monad, dyad, triad as
  fundamental number principles.)
\item
  Peirce, C. S. \emph{Collected Papers of Charles Sanders Peirce.} Vol.
  1. Harvard University Press, 1931. (§§300-353: Categories of
  Firstness, Secondness, Thirdness; §346: formal irreducibility of
  Thirdness.)
\item
  Shannon, C. E. ``A Mathematical Theory of Communication.'' \emph{Bell
  System Technical Journal} 27(3): 379-423, 1948. (Cited in
  \href{https://github.com/baardev/truevalue/raw/main/docnav/Research/papers/3_minimal-recursive-triadic-framework/3_minimal-recursive-triadic-framework.pdf}{paper
  3}: minimum bit for encoding a relation.)
\item
  Spencer-Brown, G. \emph{Laws of Form.} Allen and Unwin, 1969.
  (Distinction operator, marked/unmarked states, and re-entry as the
  minimum for self-referential dynamics.)
\item
  Von Neumann, J. ``Zur Einführung der transfiniten Zahlen.'' \emph{Acta
  Scientiarum Mathematicarum} 1: 199-208, 1923. (Construction of
  ordinals as sets: \(0 = \emptyset\), \(1 = \{0\}\), \(2 = \{0,1\}\),
  \(3 = \{0,1,2\}\).)
\end{enumerate}

\begin{center}\rule{0.5\linewidth}{0.5pt}\end{center}

\subsection{Appendix: Notation and cross-paper
reference}\label{appendix-notation-and-cross-paper-reference}

\begin{longtable}[]{@{}
  >{\raggedright\arraybackslash}p{(\linewidth - 4\tabcolsep) * \real{0.2500}}
  >{\raggedright\arraybackslash}p{(\linewidth - 4\tabcolsep) * \real{0.2812}}
  >{\raggedright\arraybackslash}p{(\linewidth - 4\tabcolsep) * \real{0.4688}}@{}}
\toprule\noalign{}
\begin{minipage}[b]{\linewidth}\raggedright
Symbol
\end{minipage} & \begin{minipage}[b]{\linewidth}\raggedright
Meaning
\end{minipage} & \begin{minipage}[b]{\linewidth}\raggedright
First defined
\end{minipage} \\
\midrule\noalign{}
\endhead
\bottomrule\noalign{}
\endlastfoot
\(N\) & Negotiation variable (running state) &
\href{https://github.com/baardev/truevalue/raw/main/docnav/Research/papers/1_recursive-tholonic-five-constants/1_recursive-tholonic-five-constants.pdf}{Paper
1}, §2 \\
\(D\) & Definition/limitation variable (bounding) &
\href{https://github.com/baardev/truevalue/raw/main/docnav/Research/papers/1_recursive-tholonic-five-constants/1_recursive-tholonic-five-constants.pdf}{Paper
1}, §2 \\
\(C\) & Contribution/integration variable (accumulating) &
\href{https://github.com/baardev/truevalue/raw/main/docnav/Research/papers/1_recursive-tholonic-five-constants/1_recursive-tholonic-five-constants.pdf}{Paper
1}, §2 \\
\(\mathbf{n}\) & Ordinal category with \(n\) objects & Mac Lane
{[}ML98{]} \\
\(K_n\) & Complete graph on \(n\) vertices & Standard graph theory \\
\(\Delta^k\) & \(k\)-simplex & Hatcher {[}Hat02{]} \\
\(S_n\) & Symmetric group on \(n\) elements & Armstrong {[}Arm88{]} \\
\(\mu(G)\) & Cyclomatic complexity of graph \(G\): \(E - V + 1\) &
Standard graph theory \\
Lemma 4.2 & Three variables are necessary for non-trivial convergence &
\href{https://github.com/baardev/truevalue/raw/main/docnav/Research/papers/3_minimal-recursive-triadic-framework/3_minimal-recursive-triadic-framework.pdf}{Paper
3}, §4.1 \\
\end{longtable}

\end{document}
